\documentclass[preview]{standalone}

\usepackage{amsmath}
\usepackage{amssymb}
\usepackage{stellar}
\usepackage{definitions}
\usepackage{bettelini}

\begin{document}

\id{probability-event-independence}
\genpage

\section{Event independence}

\begin{snippetdefinition}{independent-events-definition}{Independent events}
    Let \((\Omega,\mathcal F,\mathbb P)\) be a \probabilityspace.
    Two \event[events] \(A,B\in\mathcal F\) are
    \emph{independent} if
    \[
        \mathbb P(A\intersection B)
        =
        \mathbb P(A)\mathbb P(B).
    \]
\end{snippetdefinition}

\begin{snippetproposition}{conditional-characterization-independent-events}{Conditional characterization of independent events}
    Let \((\Omega,\mathcal F,\mathbb P)\) be a \probabilityspace and let
    \(A,B\in\mathcal F\).
    If \(\mathbb P(B)>0\), then \(A\) and \(B\) are \independentevents
    if and only if
    \[
        \mathbb P(A\mid B)=\mathbb P(A).
    \]
    If also \(\mathbb P(A)>0\), this is equivalent to
    \[
        \mathbb P(B\mid A)=\mathbb P(B).
    \]
\end{snippetproposition}

\begin{snippetproof}{conditional-characterization-independent-events-proof}{conditional-characterization-independent-events}{Conditional characterization of independent events}
    Assume \(\mathbb P(B)>0\). By definition,
    \[
        \mathbb P(A\mid B)=\mathbb P(A)
    \]
    if and only if
    \[
        \frac{\mathbb P(A\intersection B)}{\mathbb P(B)}
        =
        \mathbb P(A),
    \]
    which is equivalent to
    \[
        \mathbb P(A\intersection B)
        =
        \mathbb P(A)\mathbb P(B).
    \]
    This is precisely the definition of \independentevents.
    The statement with \(\mathbb P(B\mid A)\) follows by the same argument
    when \(\mathbb P(A)>0\).
\end{snippetproof}

\begin{snippetproposition}{disjoint-independent-events}{Disjoint independent events}
    Let \((\Omega,\mathcal F,\mathbb P)\) be a \probabilityspace and let
    \(A,B\in\mathcal F\). If \(A\intersection B=\emptyset\), then \(A\) and
    \(B\) are \independentevents if and only if
    \[
        \mathbb P(A)=0
        \quad\text{or}\quad
        \mathbb P(B)=0.
    \]
\end{snippetproposition}

\begin{snippetproof}{disjoint-independent-events-proof}{disjoint-independent-events}{Disjoint independent events}
    If \(A\intersection B=\emptyset\), then
    \[
        \mathbb P(A\intersection B)=0.
    \]
    Thus independence is equivalent to
    \[
        0=\mathbb P(A)\mathbb P(B),
    \]
    which holds if and only if \(\mathbb P(A)=0\) or \(\mathbb P(B)=0\).
\end{snippetproof}

\begin{snippetdefinition}{independent-family-events-definition}{Independent family of events}
    Let \((\Omega,\mathcal F,\mathbb P)\) be a \probabilityspace.
    The \event[events] \(A_1,\ldots,A_n\in\mathcal F\) are
    \emph{independent} if, for every non-empty
    \(J\subseteq\{1,\ldots,n\}\),
    \[
        \mathbb P\left(\bigcap_{j\in J}A_j\right)
        =
        \prod_{j\in J}\mathbb P(A_j).
    \]
\end{snippetdefinition}

\begin{snippetexample}{pairwise-independent-not-mutually-independent-events}{Pairwise independent events need not be mutually independent}
    Let \(\Omega=\{H,T\}^3\) be the \samplespace for three distinguishable fair
    coin tosses, equipped with \classicalprobability.
    Define the \event[events]
    \[
        A_1=\{\text{the first and second tosses are equal}\},
    \]
    \[
        A_2=\{\text{the first and third tosses are equal}\},
    \]
    and
    \[
        A_3=\{\text{the second and third tosses are equal}\}.
    \]
    Then
    \[
        \mathbb P(A_1)=\mathbb P(A_2)=\mathbb P(A_3)=\frac12,
    \]
    and every pairwise intersection has probability \(\frac14\). For example,
    \(A_1\intersection A_2\) is the \event that all three tosses are equal.
    Hence the \event[events] are pairwise \independentevents.
    However,
    \[
        A_1\intersection A_2\intersection A_3
        =
        \{(H,H,H),(T,T,T)\},
    \]
    so
    \[
        \mathbb P(A_1\intersection A_2\intersection A_3)
        =
        \frac14
        \neq
        \frac18
        =
        \mathbb P(A_1)\mathbb P(A_2)\mathbb P(A_3).
    \]
    Therefore \(A_1,A_2,A_3\) are not an \independentfamilyofevents.
\end{snippetexample}

\begin{snippetexample}{total-intersection-not-enough-for-independence}{The total intersection condition is not enough}
    Let \(\Omega=\{1,\ldots,8\}\), equipped with \classicalprobability.
    Define
    \[
        A=\{1,2,3,4\},\qquad
        B=\{1,5,6,7\},\qquad
        C=\{1,2,3,8\}.
    \]
    Then
    \[
        \mathbb P(A)=\mathbb P(B)=\mathbb P(C)=\frac12,
    \]
    and
    \[
        \mathbb P(A\intersection B\intersection C)
        =
        \frac18
        =
        \mathbb P(A)\mathbb P(B)\mathbb P(C).
    \]
    Nevertheless,
    \[
        A\intersection B=\{1\},
        \qquad
        \mathbb P(A\intersection B)=\frac18
        \neq
        \frac14
        =
        \mathbb P(A)\mathbb P(B).
    \]
    Thus the factorization of the intersection of all three \event[events]
    does not imply independence of the family.
\end{snippetexample}

\begin{snippettheorem}{independence-complements-characterization}{Independence and complements}
    Let \((\Omega,\mathcal F,\mathbb P)\) be a \probabilityspace and let
    \(A_1,\ldots,A_n\in\mathcal F\). The following are equivalent:
    \begin{enumerate}
        \item \(A_1,\ldots,A_n\) are an \independentfamilyofevents;
        \item for every choice \(B_i\in\{A_i,\Omega\difference A_i\}\),
            \[
                \mathbb P\left(\bigcap_{i=1}^n B_i\right)
                =
                \prod_{i=1}^n\mathbb P(B_i).
            \]
    \end{enumerate}
\end{snippettheorem}

\begin{snippetproof}{independence-complements-characterization-proof}{independence-complements-characterization}{Independence and complements}
    Suppose first that \(A_1,\ldots,A_n\) are independent.
    We show that complementing any subcollection preserves the product formula.
    It is enough to complement one \event at a time. Let
    \(C\) be the intersection of a subcollection of the remaining \event[events]
    and their complements, and suppose that
    \[
        \mathbb P(C\intersection A_k)=\mathbb P(C)\mathbb P(A_k).
    \]
    Then
    \[
        C\intersection(\Omega\difference A_k)
        =
        C\difference(C\intersection A_k),
    \]
    and therefore
    \[
        \mathbb P(C\intersection(\Omega\difference A_k))
        =
        \mathbb P(C)-\mathbb P(C\intersection A_k)
        =
        \mathbb P(C)(1-\mathbb P(A_k))
        =
        \mathbb P(C)\mathbb P(\Omega\difference A_k).
    \]
    Iterating this argument gives the product formula for every choice
    \(B_i\in\{A_i,\Omega\difference A_i\}\).

    Conversely, assume the product formula holds for every such choice.
    To prove that \(A_1,\ldots,A_n\) are independent, let
    \(\emptyset\neq J\subseteq\{1,\ldots,n\}\). Summing the product formula
    over all choices \(B_i\in\{A_i,\Omega\difference A_i\}\) for
    \(i\notin J\), and keeping \(B_j=A_j\) for \(j\in J\), gives
    \[
        \mathbb P\left(\bigcap_{j\in J}A_j\right)
        =
        \prod_{j\in J}\mathbb P(A_j)
        \prod_{i\notin J}
        \bigl(\mathbb P(A_i)+\mathbb P(\Omega\difference A_i)\bigr).
    \]
    Since each factor in the last product is \(1\), the desired factorization
    for the subfamily indexed by \(J\) follows.
\end{snippetproof}

\begin{snippetcorollary}{independence-preserved-by-complements}{Independence is preserved by complements}
    Let \((\Omega,\mathcal F,\mathbb P)\) be a \probabilityspace and let
    \(A_1,\ldots,A_n\in\mathcal F\) be independent \event[events].
    If \(B_i\in\{A_i,\Omega\difference A_i\}\) for every \(i\), then
    \(B_1,\ldots,B_n\) are also independent \event[events].
\end{snippetcorollary}

\begin{snippetproof}{independence-preserved-by-complements-proof}{independence-preserved-by-complements}{Independence is preserved by complements}
    This is exactly the forward implication of the characterization of
    independence by complements.
\end{snippetproof}

\end{document}
